\documentclass{exam}
\usepackage{amsmath, amssymb}

\pagestyle{headandfoot}
\firstpageheadrule
\runningheadrule
\firstpageheader{Prof. Tahvildar-Zadeh \\ Linear Algebra}{Homework 9}{Aadhithya Saravanan}
\runningheader{Linear Algebra \\ Homework 9}{}{Saravanan}
\firstpagefooter{}{}{}
\runningfooter{}{\thepage}{}

\printanswers

\begin{document}

\underline{Exercise 5.4.2}
\newline
\begin{parts}
    \part To check whether $W$ is a $T$-invariant subspace of $V$, we can check whether each vector in a basis for $W$ remains in $W$ after $T$ is applied. Let $\beta=\{1,x,x^2\}$ be a basis for $W$. Then,
    \[T(1)=0\in\text{span}(\beta)\]
    \[T(x)=1\in\text{span}(\beta)\]
    \[T(x^2)=2x\in\text{span}(\beta)\]
    So, $T(v)\in W$ for all $v\in W$, so $T(W)\subseteq W$. Therefore, $W$ is a $T$-invariant subspace of $V$.
    \part To check whether $W$ is a $T$-invariant subspace of $V$, we can check whether each vector in a basis for $W$ remains in $W$ after $T$ is applied. Let $\beta=\{1,x,x^2\}$ be a basis for $W$. Then,
    \[T(1)=x\in\text{span}(\beta)\]
    \[T(x)=x^2\in\text{span}(\beta)\]
    \[T(x^2)=x^3\notin\text{span}(\beta)\]
    So, $T(v)\notin W$ for all $v\in W$, so $T(W)\nsubseteq W$. Therefore, $W$ is not a $T$-invariant subspace of $V$.
    \part To check whether $W$ is a $T$-invariant subspace of $V$, we can check whether each vector in a basis for $W$ remains in $W$ after $T$ is applied. Let $\beta=\left\{
    \begin{bmatrix}
        1 \\
        1 \\
        1
    \end{bmatrix}
    \right\}$ be a basis for $W$. Then,
    \[T\left(
    \begin{bmatrix}
        1 \\
        1 \\
        1
    \end{bmatrix}
    \right)=
    \begin{bmatrix}
        3 \\
        3 \\
        3
    \end{bmatrix}
    \in\text{span}(\beta)\]
    So, $T(v)\in W$ for all $v\in W$, so $T(W)\subseteq W$. Therefore, $W$ is a $T$-invariant subspace of $V$.
    \newline
\end{parts}

\underline{Exercise 5.4.3}
\newline
\begin{parts}
    \part Let $W=\{0\}$. Since $T$ is linear, $T(0)=0$. So, for all $v\in W$, $T(v)\in W$. Then, $T(W)\subseteq W$. Therefore, $W$ is a $T$-invariant subspace of $V$.
    \newline
    \newline
    Let $\beta$ be a basis for $V$. Since $T$ is a linear operator on $V$, it maps vectors in $V$ to vectors in $V$. So, every vector in $\beta$ maps to a vector in $V$. Then, for all $v\in V$, $T(v)\in V$. Then, $T(V)\subseteq V$. Therefore, $V$ is a $T$-invariant subspace of $V$.
    \part Let $W=N(T)=\{v\in V\mid T(v)=0\}$. Since $T$ is linear, $T(0)=0$. So, $0\in N(T)$. For all $v\in W$, $T(v)=0\in N(T)=W$. Then, $T(W)\subseteq W$. Therefore, $W$ is a $T$-invariant subspace of $V$.
    \newline
    \newline
    Let $W=R(T)=\{T(v)\mid v\in V\}$. For all $v\in V$, $T(v)\in R(T)=W$, by definition of $R(T)$. Since $W$ is a subspace of $V$, for all $v\in W$, $v\in V$. So, by the definition of $R(T)$, $T(v)\in R(T)=W$. Then, $T(W)\subseteq W$. Therefore, $W$ is a $T$-invariant subspace of $V$.
    \part Let $W=E_\lambda=\{v\in V\mid T(v)=\lambda v\}$. For all $v\in W$, $T(v)=\lambda v\in E_\lambda=W$. Then, $T(W)\subseteq W$. Therefore, $W$ is a $T$-invariant subspace of $V$.
    \newline
\end{parts}

\underline{Exercise 5.4.6}
\newline
\begin{parts}
    \part We want to find $\beta$, an ordered basis for the $T$-cyclic subspace generated by $e_1$. Let $\beta_1=e_1$.
    \[T(e_1)=T\left(
    \begin{bmatrix}
        1 \\
        0 \\
        0 \\
        0
    \end{bmatrix}
    \right)=
    \begin{bmatrix}
        1 \\
        0 \\
        1 \\
        1
    \end{bmatrix}=\beta_2\]
    \[T(\beta_2)=T\left(
    \begin{bmatrix}
        1 \\
        0 \\
        1 \\
        1
    \end{bmatrix}
    \right)=
    \begin{bmatrix}
        1 \\
        -1 \\
        2 \\
        2
    \end{bmatrix}=\beta_3\]
    \[T(\beta_3)=T\left(
    \begin{bmatrix}
        1 \\
        -1 \\
        2 \\
        2
    \end{bmatrix}
    \right)=
    \begin{bmatrix}
        0 \\
        -3 \\
        3 \\
        3
    \end{bmatrix}=-3\beta_2+3\beta_3\]
    So, $\beta=\{\beta_1,\beta_2,\beta_3\}$ is an ordered basis for the $T$-cyclic subspace generated by $e_1$.
    \part We want to find $\beta$, an ordered basis for the $T$-cyclic subspace generated by $x^3$. Let $\beta_1=x^3$.
    \[T(x^3)=6x=\beta_2\]
    \[T(\beta_2)=0=0\beta_1\]
    So, $\beta=\{\beta_1,\beta_2\}$ is an ordered basis for the $T$-cyclic subspace generated by $x^3$.
    \newline
\end{parts}

\underline{Exercise 5.4.9}
\newline
\begin{parts}
    \part \[\beta=\{\beta_1,\beta_2,\beta_3\}=\left\{
    \begin{bmatrix}
        1 \\
        0 \\
        0 \\
        0
    \end{bmatrix},
    \begin{bmatrix}
        1 \\
        0 \\
        1 \\
        1
    \end{bmatrix},
    \begin{bmatrix}
        1 \\
        -1 \\
        2 \\
        2
    \end{bmatrix}
    \right\}
    \]
    We know that $T^3(e_1)=-3\beta_2+3\beta_3$. So,
    \[0e_1+3T(e_1)-3T^2(e_1)+T^3(e_1)=0\]
    By Theorem 5.21.b, $T_W(t)=(-1)^3(0+3t-3t^2+t^3)=-t^3+3t^2-3t$.
    \newline
    \newline
    Since $T(\beta_1)=\beta_2$, $T(\beta_2)=\beta_3$, and $T(\beta_3)=-3\beta_2+3\beta_3$,
    \[[T_W]_\beta=
    \begin{bmatrix}
        0 & 0 & 0 \\
        1 & 0 & -3 \\
        0 & 1 & 3
    \end{bmatrix}\]
    So, $T_W(t)=\det([T_W]\beta-tI)=-t((-t)(3-t)-(-3))=-t(t^2-3t+3)=-t^3+3t^2-3t$.
    \part \[\beta=\{\beta_1,\beta_2\}=\{x^3,6x\}\]
    We know that $T^2(x^3)=0=0\beta_1$. So,
    \[0\beta_1+0T(\beta_1)+T^2(\beta_1)=0\]
    By Theorem 5.21.b, $T_W(t)=(-1)^2(0+0t+t^2)=t^2$
    \newline
    \newline
    Since $T(\beta_1)=\beta_2$ and $T(\beta_2)=0$,
    \[[T_W]_\beta=
    \begin{bmatrix}
        0 & 0 \\
        1 & 0
    \end{bmatrix}\]
    So, $T_W(t)=\det([T_W]\beta-tI)=-t(-t)-0(1)=t^2$.
    \newline
\end{parts}

\underline{Exercise 5.4.10}
\newline
\begin{parts}
    \part Let $\beta$ be the standard ordered basis for $\mathbb{R}^4$.
    \[T(e_1)=T\left(
    \begin{bmatrix}
        1 \\
        0 \\
        0 \\
        0 
    \end{bmatrix}
    \right)=
    \begin{bmatrix}
        1 \\
        0 \\
        1 \\
        1 
    \end{bmatrix}
    =e_1+e_3+e_4
    \]
    \[T(e_2)=T\left(
    \begin{bmatrix}
        0 \\
        1 \\
        0 \\
        0 
    \end{bmatrix}
    \right)=
    \begin{bmatrix}
        1 \\
        1 \\
        0 \\
        0 
    \end{bmatrix}
    =e_1+e_2
    \]
    \[T(e_3)=T\left(
    \begin{bmatrix}
        0 \\
        0 \\
        1 \\
        0 
    \end{bmatrix}
    \right)=
    \begin{bmatrix}
        0 \\
        -1 \\
        1 \\
        0 
    \end{bmatrix}
    =-e_2+e_3
    \]
    \[T(e_4)=T\left(
    \begin{bmatrix}
        0 \\
        0 \\
        0 \\
        1 
    \end{bmatrix}
    \right)=
    \begin{bmatrix}
        0 \\
        0 \\
        0 \\
        1 
    \end{bmatrix}
    =e_4
    \]
    So,
    \[
    [T]_\beta=
    \begin{bmatrix}
        1 & 1 & 0 & 0 \\
        0 & 1 & -1 & 0 \\
        1 & 0 & 1 & 0 \\
        1 & 0 & 0 & 1
    \end{bmatrix}
    \]
    Then, $f(t)=\det([T]_\beta-tI)=(1-t)((1-t)((1-t)(1-t))+1(-1))=(1-t)(-t^3+3t^2-3t)=(1-t)T_W(t)$. So, $T_W(t)$ divides $f(t)$.
    \part Let $\beta=\{1,x,x^2,x^3\}$ be an ordered basis for $P_3(\mathbb{R})$.
    \[T(1)=0\]
    \[T(x)=0\]
    \[T(x^2)=2=2\beta_1\]
    \[T(x^3)=6x=6\beta_2\]
    So,
    \[
    [T]_\beta=
    \begin{bmatrix}
        0 & 0 & 2 & 0 \\
        0 & 0 & 0 & 6 \\
        0 & 0 & 0 & 0 \\
        0 & 0 & 0 & 0
    \end{bmatrix}
    \]
    Then, $f(t)=\det([T]_\beta-tI)=t^4=t^2(t^2)=t^2T_W(t)$. So, $T_W(t)$ divides $f(t)$.
    \newline
\end{parts}

\underline{Exercise 5.4.18}
\newline
\begin{parts}
    \part First, we will prove the forward direction. Assume $A$ is invertible. We want to show that $a_0$ is not equal to 0. Let $T$ be the linear transformation associated with $A$, such that $T(v)=Av$. Since $A$ is invertible, $T$ is invertible. $T$ is also one-to-one. By Theorem 2.5, $\text{rank}(T)=\dim(V)$, where $V$ is a vector space with $n$ dimensions. By the Dimension Theorem, $\dim(N(T))+\dim(R(T))=dim(V)$. Since, $\dim(R(T))=\text{rank}(T)$, $\dim(N(T))=0$. Therefore, $T(v)=0$ if and only if $v=0$. Assume $a_0=0$. Then, $f(t)$ can be factored into $t((-1)^nt^{n-1}+a_{n-1}t^{n-2}+\dots+a_1)$. Then, 0 is an eigenvalue of $T$. Then, there exists a nonzero vector $v$ such that $T(v)=0$. However, this is a contradiction, since $T(v)=0$ if and only if $v=0$. So, $a_0\ne0$.
    \newline
    \newline
    Now, we will prove the backward direction. Assume $a_0$ is not equal to 0. We want to show that $A$ is invertible. Let $T$ be the linear transformation associated with $A$, such that $T(v)=Av$. Since $a_0$ is not 0, a $t$ cannot be factored out of $f(t)$. So, 0 is not an eigenvalue of $T$. Then, there does not exist a nonzero vector $v$ such that $T(v)=0$. So, $N(T)=\{0\}$. By Theorem 2.4, $T$ is one-to-one. By Theorem 2.5, $T$ is onto. Since $T$ is bijective, $T$ is invertible. Then, $A$ must be invertible, as the associated matrix of a linear transformation is invertible if and only if the linear transformation is invertible. Therefore, $A$ is invertible if and only if $a_0\ne0$.
    \part By the Cayley-Hamilton Theorem for Matrices, $f(A)=O$. So,
    \[O=(-1)^nA^n+a_{n-1}A^{n-1}+\dots+a_1A+a_0I\]
    \[-a_0I=(-1)^nA^n+a_{n-1}A^{n-1}+\dots+a_1A\]
    \[A^{-1}(-a_0I)=A^{-1}((-1)^nA^n+a_{n-1}A^{n-1}+\dots+a_1A)\]
    \[-a_0A^{-1}=(-1)^nA^{n-1}+a_{n-1}A^{n-2}+\dots+a_1I\]
    \[A^{-1}=-\frac{1}{a_0}[(-1)^nA^{n-1}+a_{n-1}A^{n-2}+\dots+a_1I]\]
    \part $f(t)=\det(A-tI)=(1-t)(2-t)(-1-t)=-t^3+2t^2+t-2$. So,
    \[A^{-1}=-\frac{1}{-2}[(-1)^3A^2+2A+I]\]
    \[A^{-1}=\frac{1}{2}[-A^2+2A+I]\]
    \[A^{-1}=\frac{1}{2}\left[-
    \begin{bmatrix}
        1 & 2 & 1 \\
        0 & 2 & 3 \\
        0 & 0 & -1
    \end{bmatrix}^2+2
    \begin{bmatrix}
        1 & 2 & 1 \\
        0 & 2 & 3 \\
        0 & 0 & -1
    \end{bmatrix}+
    \begin{bmatrix}
        1 & 0 & 0 \\
        0 & 1 & 0 \\
        0 & 0 & 1
    \end{bmatrix}
    \right]\]
    \[A^{-1}=\frac{1}{2}\left[
    \begin{bmatrix}
        -1 & -6 & -6 \\
        0 & -4 & -3 \\
        0 & 0 & -1
    \end{bmatrix}+
    \begin{bmatrix}
        2 & 4 & 2 \\
        0 & 4 & 6 \\
        0 & 0 & -2
    \end{bmatrix}+
    \begin{bmatrix}
        1 & 0 & 0 \\
        0 & 1 & 0 \\
        0 & 0 & 1
    \end{bmatrix}
    \right]\]
    \[A^{-1}=\frac{1}{2}
    \begin{bmatrix}
        2 & -2 & -4 \\
        0 & 1 & 3 \\
        0 & 0 & -2
    \end{bmatrix}
    \]
    \[A^{-1}=
    \begin{bmatrix}
        1 & -1 & -2 \\
        0 & \frac{1}{2} & \frac{3}{2} \\
        0 & 0 & -1
    \end{bmatrix}
    \]
    \newline
\end{parts}

\end{document}